(1) SH位移满足 \(\rho_j u_{tt}=\mu_j(u_{xx}+u_{zz})\)，代入行波形式得
\(U_j''+(\omega^2/c_j^2-\beta^2)U_j=0\)。
要在覆层振荡、在基底衰减，须
\[
q^2=\omega^2/c_1^2-\beta^2>0,\quad
\alpha^2=\beta^2-\omega^2/c_2^2>0,\quad
c_1<\omega/\beta<c_2.
\]
自由表面 \(U_1'(0)=0\) 选 \(U_1=A\cos(qz)\)，基底束缚解 \(U_2=D e^{-\alpha(z-h)}\)，舍弃无穷远增长分支。

(2) \(z=h\) 位移连续给 \(D=A\cos(qh)\)，剪应力连续给
\(-\mu_1Aq\sin(qh)=-\mu_2\alpha D\)，因此
\[
\boxed{\tan(qh)=\frac{\mu_2\alpha}{\mu_1q}}.
\]
两边无量纲，模态同时依赖剪切波速与剪切模量比。若错误地用纵波法向应力替代SH剪应力，所得色散式将失去正确的横向极化边界条件。

(3) 第 \(n\ge1\) 阶模截止时 \(\alpha\to0^+\)，相速度趋 \(c_2\)，此时 \(\tan(qh)\to0\)、\(qh=n\pi\)。故
\[
\boxed{\omega_{n,\rm cut}=\frac{n\pi}{h\sqrt{c_1^{-2}-c_2^{-2}}}},\quad n\ge1.
\]
基模 \(n=0\) 无非零截止频率。低频时 \(q h\to0\)、\(\alpha\to0\)，相速度趋 \(c_2^-\)；高频基模 \(q h\to\pi/2^-\)，\(\beta\simeq\omega/c_1\)，相速度趋 \(c_1^+\)。

(4) 单位水平面积的周期平均能量可写
\[
\mathcal E=\frac14\int_0^\infty[\rho(z)\omega^2|U|^2+
\mu(z)(\beta^2|U|^2+|U'|^2)]\,dz.
\]
基底位移平方积分为 \(|D|^2/(2\alpha)\)，故衰减长度 \(1/\alpha\)。截止附近 \(\alpha\to0\) 时基底能量延伸至远处，局域性减弱；\(c_1\to c_2\) 使允许的相速度窗口闭合且高模截止频率发散，均匀半空间不能支持同类Love束缚波。
